\section{齐：海边的国家，先学会怎样召集天下}

齐在东方，背后是海，面前是中原。它有盐、鱼、交通和相对富庶的腹地，却没有资格天然成为霸主。周王室东迁之后，谁能在边患来临时出兵、在诸侯相持时召集会面，谁才有机会把本国的力量说成“天下的事”。齐桓公和管仲的时代，正是在这个空缺里出现。

\eventref{SA-0685-QIHUAN}{春秋·齐桓公即位}

齐襄公被杀后，公子小白与纠争位；小白先入齐即位，是为齐桓公，原先依附公子纠的管仲后来也进入齐国执政集团。这一转折的年次和结果可由《左传·庄公九年》的编年脉络把握；《国语·齐语》把齐桓、管仲置入一套“修政以服诸侯”的国别政治记忆；《史记·齐太公世家》又将它收束为霸业的开端。三种材料都不是宫门内的逐时记录。可以较有把握地说的，是一次继承危机之后，齐国重新取得了能够调度内部与对外交涉的中心；至于每一项整顿怎样制定、由谁独断，传世材料不足以列成一张行政清单。

齐的优势在于，它把内部整合和对外会盟接了起来。一个诸侯国若只谈本国富强，会引起别国警惕；若能在救援、迁徙与会盟中让别国看见可得的帮助，又把出兵说成“尊王攘夷”，自身利益便能取得公共秩序的语言。齐桓公的霸权从来不是一个人站在高处发号施令，而是许多国家在周王室无力之后，按各自的险急和利害，暂时愿意让齐来组织局面。它也不是一张自始固定的盟国名册：留存的经传更多保存了某次救援、会师或结盟的节点，而非一份逐年更新的从属簿。

\eventanchor{SA-0664-QI-YAN}{春秋·齐援燕伐山戎}

约前664年，北方的燕受到山戎压力。《左传·庄公三十年》以“伐山戎，救燕”留下齐军北出的最小叙事；《史记·齐太公世家》将远征延展为更完整的战功线，所涉孤竹、行军及归还土地的细节须留在汉代整理的叙事层次中。即使不把后出的行军故事当作地图，齐能为燕出兵这件事本身，已让北方诸侯看到齐的兵力可以越出本土，介入一个本不由齐直接统治的危机。

\eventanchor{SA-0659-QI-XING}{春秋·齐迁邢于夷仪}

五年后，狄人伐邢。《春秋·僖公元年》记录了狄人的进攻以及齐、宋、曹的救援行动；《左传》才把邢的迁徙安放到夷仪。城邑被迫挪动，意味着受援者面对的不只是一次战斗，而是原有居处与安全条件的破裂；但迁徙规模、具体工程和齐军的全部行动，不能由几句经传补成现代意义的安置档案。齐的作用同样不宜写成施恩：救燕与援邢给齐带来威望和关系，也让燕、邢等国在下一次危机中有理由期待齐的召集。援助、声望与会盟由此互相增益，却仍要靠每次出兵和每个诸侯重新兑现。

这条北方关系网，是齐军南下的使能条件，不是南下的唯一原因。齐并未因此取得支配诸侯内政的权利；它取得的是在共同危险面前提出议程的资格。到前656年，齐能把若干诸侯的兵力带向楚境，正是这种资格在更大范围内接受检验的时刻。

\begin{event}{公元前656年}{召陵之盟}{可靠纪年}{召陵一带}\eventanchor{SA-0656-SHAOLING}{春秋·召陵之盟}
《春秋·僖公四年》把会师、侵蔡、伐楚、驻军与“楚屈完来盟于师，盟于召陵”排成短促的编年句子，足以确认齐主导的联盟到了楚的北面，也足以确认双方以盟而非吞并收束。它并不告诉我们每一支军队如何行进，或楚人为何在这一刻同意交涉。《左传》补入齐方以不贡苞茅责问楚使的场景，因而把问题置于周礼贡赋的框架：谁能要求楚承认这套语言，谁又能代表它说话。《国语·齐语》所保存的齐国政治记忆，和《史记》后来写成的齐霸长线，都使这次交涉更像齐桓霸业的高峰；它们可以说明后人怎样理解召陵，不能反过来替代经文所未记的谈判记录。
\end{event}

\begin{textwindow}[水滨的反问]
“昭王之不复，君其问诸水滨。”

出处为《左传·僖公四年》：叙事让楚使以周昭王南征未归反问齐军，正嵌在双方争执贡茅的场景里。它把一场外交交锋写得锋利，却是后出编年叙事所塑造的使者言辞；可借以看见周礼权威如何在楚人面前受到质问，不能把它当作会谈记录。
\end{textwindow}

楚没有被齐吞并，也没有因此变成齐的属国。屈完来盟这一结果，使齐获得了能够把南方大国拉入自己主持的仪式与交涉之中的声望；楚则没有放弃其独立的军政位置。霸权在这里不是消灭对手，而是迫使对手在自己提出的语言中留下一个位置。它的边界也正在这里：一旦齐不能继续把兵力、盟友和周室名分带到现场，召陵的语言本身不会自动约束楚国。

\begin{event}{公元前651年}{葵丘之会}{可靠纪年}{葵丘}\eventanchor{SA-0651-KUIQIU}{春秋·葵丘之会}
《春秋·僖公九年》只郑重记下“齐侯盟诸侯于葵丘”。《左传》赋予盟约更具体的禁约与礼仪意义，《史记》又把周室承认、齐桓受命的故事编入一条更完整的霸业叙事。经文能够稳固地证明会盟及其年次；盟辞的全貌、周室使者的每一步仪节，以及齐桓当时的意图，则分别属于后出的叙事与汉代整合，应当和经年骨架分开阅读。
\end{event}

即使如此，葵丘仍是齐霸最像制度的一刻。盟约把彼此不尽相同的国家暂时放到同一套禁约与礼仪里，齐桓公则像是替周王室接住了协调诸侯的责任。然而，留存材料所见的是一次次会盟、救援和共同出兵，并非一个能常设征收、裁决、任命或处理继承的齐国机构。它靠的仍是齐桓公、管仲、齐国资源和诸侯的即时判断；盟员可以从齐的组织中获得安全和声望，也始终保有按本国危机转向的余地。

\begin{event}{公元前643年}{齐桓公死，诸子争位}{可靠纪年}{齐都}
\eventanchor{SA-0643-QIHUAN-DEATH}{春秋·齐桓公死与诸子争位}
《春秋·僖公十七年》保存了齐桓公卒的年次；《左传》随后把易牙、寺人貂、无亏和诸公子的竞争组织成宫廷控制权崩解的叙事。《史记·齐太公世家》又将齐桓末年写得更具兴亡转折的色彩。病中隔绝、遗体久未下葬等细节因此不宜当作无人能疑的现场报告；各层材料一致指出的，是君位一空，齐国内部已不能以一套被普遍接受的次序完成继承。
\end{event}

\eventanchor{SA-0642-QI-XIAOGONG}{春秋·齐孝公即位}

到前642年，无亏被杀，公子昭即位，是为齐孝公。《春秋·僖公十八年》提供这一年次的编年落点；无亏的身份、争位的先后和各方兵力如何进入都城，须合读《左传》与《史记》的不同组织，仍有不能一一复原之处。孝公取得君位，并不等于齐立刻取回齐桓时的诸侯号令。先前能召集天下的中心，首先没能替自己家门提供一条各方都承认的继嗣道路；诸侯也不会因一纸旧盟约而停止重新计算利害。

\begin{turningpoint}[霸权的短命，不是因为它毫无用处]
齐霸解决了周王室无力召集诸侯的难题，使会盟、救援、边患和外交有了临时的中心。它的制度得失也由此清楚：以国家能力衡量，齐把远方的求援、军队的动员和周室名分连成了可用的协调机制；代价是每次协调仍要由齐及参与诸侯承担征发、行军与承诺的成本。它的失败不在于礼仪无用，而在于礼仪不能替代可继承的制度：联盟成员不是齐的属县，盟约也不能在齐国内乱时自行运转。春秋的霸权因此既是一种创造，也是一种无法稳定传给下一代的秩序。
\end{turningpoint}

\begin{sourcenote}[齐桓、管仲与传世叙事的层次]
本节先以《春秋》的鲁国经年固定前659年、前656年、前651年以及齐桓死后的关键节点；它只能给出会、伐、盟、卒等极简骨架。《左传》提供北伐、迁邢、召陵交涉和继嗣冲突的连贯叙事，却不是会议录或宫廷档案。《国语·齐语》保存的是把齐桓与管仲塑为政治典范的国别说辞传统，适合追问合法性语言，不能逐年坐实。《史记·齐太公世家》则以汉代史家的眼光把这些材料编成一条霸业兴亡线，最能显示后世如何理解齐桓，也最需避免倒灌为春秋现场。至于《管子》，它是战国至汉代逐层形成的复合文本；其中精密的经济理论不能直接当作前685年的政策清单。齐国确实比许多诸侯更能把资源、军政与外交组织起来，但“每一条改革都由管仲亲定”的整齐故事，属于后世的加工。
\end{sourcenote}
